\begin{abstract}
% We present PhyPrompt, a lightweight, model-agnostic front end that automatically refines text-to-video prompts for improved physical realism. Our approach employs a two-stage methodology.
% First, we fine-tune a large language model on a curated, physics-focused Chain-of-Thought dataset, enabling it to integrate core physics principles such as object motion and force interactions into user prompts while preserving their original intent.
% Second, we apply Group Relative Policy Optimization with a dynamic reward schedule that initially prioritizes semantic fidelity and then progressively shifts emphasis toward commonsense physics, driving the model to correct physical inconsistencies across diverse video generation models.
% Through zero-shot evaluation on the VideoPhy2 benchmark across  a few diffusion generators (Lavie, VideoCrafter2, CogVideoX), PhyPrompt demonstrates a remarkable 16.8\% improvement in combined semantic and physical success compared to baseline prompts.
% These results show that an RL-trained LLM can steer video synthesis toward adherence to real-world physical laws, without the need for manual prompt engineering.

% State-of-the-art text-to-video (T2V) generators frequently violate physical laws despite high visual quality. We demonstrate that this limitation stems not from model capability but from insufficient physical constraints in input prompts: manually adding physics details reliably produces physically plausible videos. However, manual prompt engineering requires expertise and does not scale.
% We present PhyPrompt, a reinforcement learning framework that automatically refines prompts to elicit physically realistic video generation. Our key innovation is a dynamic reward curriculum that resolves the fundamental conflict between semantic fidelity and physical realism. Remarkably, this curriculum achieves synergistic optimization: performance on both objectives simultaneously exceeds what either objective achieves when optimized alone (Semantic Adherence: 47.8\% vs. 45.6\% from SA-only; Physical Commonsense: 66.8\% vs. 65.2\% from PC-only). This demonstrates that our method discovers compositional prompt structures beyond conventional multi-objective trade-offs.
% PhyPrompt employs a lightweight 7B parameter model that outperforms GPT-4o (+3.8\% joint performance) and DeepSeek-V3 (+2.2\%, despite 100$\times$ more parameters), achieving 40.8\% combined semantic-physical success on VideoPhy2 benchmarks. Our approach transfers zero-shot across diverse T2V architectures (Lavie, VideoCrafter2, CogVideoX-5B), delivering up to 16.8\% improvement over baseline prompts without generator-specific fine-tuning. These results establish that domain-specialized reinforcement learning with compositional curricula can surpass general-purpose scaling for physics-aware generation.
% State-of-the-art text-to-video (T2V) generators frequently violate physical laws despite high visual quality. We show this stems from insufficient physical constraints in prompts rather than model limitations: manually adding physics details reliably produces physically plausible videos, but requires expertise and does not scale.
% We present PhyPrompt, a reinforcement learning framework that automatically refines prompts for physically realistic generation. Our dynamic reward curriculum achieves synergistic optimization: PhyPrompt-7B reaches 40.8\% joint success on VideoPhy2 (8.6 percentage points gain), improving physical commonsense by 11 percentage points (55.8\% to 66.8\%) while simultaneously increasing semantic adherence by 4.4pp (43.4\% to 47.8\%). Remarkably, our curriculum exceeds single-objective training on both metrics, demonstrating compositional prompt discovery beyond conventional multi-objective trade-offs.
% PhyPrompt outperforms GPT-4o (+3.8\% joint) and DeepSeek-V3 (+2.2\%, 100$\times$ larger) using only 7B parameters. The approach transfers zero-shot across diverse T2V architectures (Lavie, VideoCrafter2, CogVideoX-5B) with up to 16.8\% improvement, establishing that domain-specialized reinforcement learning with compositional curricula surpasses general-purpose scaling for physics-aware generation.
尽管视觉质量已经很高，当前最先进的文本到视频（T2V）生成器仍然经常违反物理定律。我们指出，这一问题主要来自提示词中的物理约束不足，而不是模型能力本身：手工加入物理细节可以稳定地产生物理合理的视频，但这种做法需要专业知识，且难以规模化。
我们提出 PhyPrompt，这是一个两阶段强化学习框架，用于自动精炼提示词以诱导物理真实的视频生成。首先，我们在以物理为中心的思维链数据集上微调大语言模型，使其在保留用户意图的同时纳入物体运动、力的相互作用等原则。其次，我们采用组相对策略优化，并设计动态奖励课程：训练初期优先强调语义保真度，随后逐步转向物理常识。该课程带来协同优化：PhyPrompt-7B 在 VideoPhy2 上达到 40.8\% 的联合成功率（提升 8.6 个百分点），将物理常识从 55.8\% 提升到 66.8\%，同时将语义遵循度从 43.4\% 提升到 47.8\%。值得注意的是，该课程在两个指标上都超过单目标训练，说明它能够发现超越传统多目标权衡的组合式提示结构。
仅使用 7B 参数，PhyPrompt 就优于 GPT-4o（联合指标 +3.8\%）和 DeepSeek-V3（联合指标 +2.2\%，后者规模约为 100$\times$）。该方法还能在多种 T2V 架构（Lavie、VideoCrafter2、CogVideoX-5B）之间零样本迁移，最高带来 16.8\% 的改进，表明带有组合课程的领域专用强化学习可以在物理感知生成上超过通用规模扩展。
\end{abstract}
